%% kapa: exam=2 round=23 year=2012 period=3 subject=compensation_law subject_ko=감정평가 및 보상법규 seq=2 number=2 points=30
%% source: Q-Net 감정평가사 기출문제 (cst00309) artlSeq=5056014 file=raw/1008/5056014/3교시(감정평가 및 보상법규).hwp
%% conversion: kordoc 4.24.51 -> tools/gen2.py (auto); review=proofread-sonnet; flags=none
\begin{longquestion}{【문제 2】}
한국수자원공사는 「한국수자원공사법」 제9조 및 제10조에 근거하여 수도권(首都圈) 광역상수도사업 실시계획을 수립하여 국토해양부장관의 승인을 얻은 후, 1필지인 甲의 토지 8,000m\textsuperscript{2} 중 6,530m\textsuperscript{2}를 협의취득하였다. 협의취득 후 甲의 잔여지는 A지역 495m\textsuperscript{2}, B지역 490m\textsuperscript{2}, 그리고 C지역 485m\textsuperscript{2}로 산재(散在)하고 있다. (30점) \par
1) 甲은 위 잔여지의 토지가격의 감소를 이유로 손실보상을 청구하려고 한다. 이 경우 잔여지의 가격감소에 대한 甲의 권리구제방법을 설명하시오. (15점) \par
2) 호텔을 건립하기 위해 부지를 조성하고 있던 甲은 자신의 잔여지를 더 이상 종래의 사용목적대로 사용할 수 없게 되자 사업시행자와 매수에 관한 협의를 하였으나, 협의가 성립되지 아니하였다. 이에 甲은 관할 토지수용위원회에 잔여지의 수용을 청구하였지만, 관할 토지수용위원회는 이를 받아들이지 않았다. 이 경우 잔여지수용청구의 요건과 甲이 제기할 수 있는 행정소송의 형식을 설명하시오. (15점) \par
\end{longquestion}
